\documentclass[10pt,a4paper]{amsart}
\usepackage[margin=19mm]{geometry}
\usepackage{amsmath,amssymb,mathtools}
\usepackage[scr=rsfs,cal=euler]{mathalfa}
\usepackage{xfrac}
\usepackage[unicode,hidelinks,bookmarks=false]{hyperref}
\newcommand{\ie}{i.e.\ }
\newcommand{\cf}{cf.\ }
\newcommand{\resp}{respectively\ }
\newcommand{\Subsection}{\subsection}
\providecommand{\nobreakdash}{\nobreak\hbox{-}}
\setlength{\emergencystretch}{2em}
\multlinegap0pt
\usepackage{bm}
\usepackage{bbm}
\usepackage{dsfont}
\usepackage{xspace}
\usepackage{cancel}
\usepackage{mathtools}
\usepackage{amsthm}
\makeatletter
\@ifundefined{theorem}{\newtheorem{theorem}{Theorem}}{}
\@ifundefined{lemma}{\newtheorem{lemma}[theorem]{Lemma}}{}
\makeatother
\makeatletter
\newcommand{\R}{\mathbb{R}}
\expandafter\ifx\csname defproof\endcsname\relax
\newenvironment{defproof}[1]{\par\noindent{\bf Proof~#1:}}{\qed \par\medskip\noindent}
\fi
\expandafter\ifx\csname proof\endcsname\relax
\newenvironment{proof}{\par\noindent{\bf Proof:}}{\qed \par\medskip\noindent}
\fi
\expandafter\ifx\csname proofof\endcsname\relax
\newenvironment{proofof}[1]{\par\noindent{\bf Proof of #1:}}{\qed \par\medskip\noindent}
\fi
\makeatother
\title[On the quality of randomized approximations of Tukey's depth]{On the quality of randomized approximations of Tukey's depth}
\author{Simon Briend, G\'abor Lugosi, Roberto Imbuzeiro Oliveira}
\date{Source: arXiv:2309.05657v3 (CC BY 4.0)}
\begin{document}
\maketitle

\noindent\emph{Source and licence.} On the quality of randomized approximations of Tukey's depth. Version arXiv:2309.05657v3 (CC BY 4.0), distributed under the Creative Commons Attribution 4.0 International licence, as recorded in the arXiv licence metadata for this exact version and shown on the abstract page https://arxiv.org/abs/2309.05657v3. Full rights notice: licenses/arxiv-2309.05657-CC-BY-4.0.txt.\par\smallskip

\noindent\textit{Editorial scope.} Counted unit: the Lemma whose source label is \texttt{density lower bound} in the article (source0.tex lines 1203--1208, source label \texttt{density lower bound}), delivered together with the article's own complete original proof of it (source0.tex lines 1211--1277, one \texttt{proof} environment of 67 lines). No mathematics is written, repaired or re-derived: statement and proof are the article's own text line for line. Required context retained uncounted: none. Every definition, display and label of those lines is retained unchanged. Omitted from this package and not redistributed: 1--1202, 1209--1210, 1278--1529. Nothing inside a retained excerpt is omitted or altered: every retained body is delivered exactly as the article's lines are written, so no omission receipt of any kind accompanies this package. Citations: the retained excerpts carry no citation command at all, so no bibliography entry of the article is transcribed and no reference is redistributed. Reference adaptation: none was needed and none was made; every reference inside the delivered unit resolves inside it, so no unresolved reference or citation is produced. Printed slips kept literal: no printed word, symbol, hypothesis or number of the counted statement or of its proof was altered. Layout and class: the article's own class is \texttt{article} with packages including inputenc, kpfonts, titlesec, pgf, tikz, hyperref, graphicx, booktabs, amsfonts, amsmath, amssymb, tcolorbox, bbm, color; the wrapper is the lane's pinned \texttt{amsart} editorial wrapper at 10pt on A4, which restates the theorem-like environments the retained text needs and adds the article's own macro definitions that the retained text needs (\texttt{\textbackslash R}), with extra wrapper packages (bm, bbm, dsfont, xspace, cancel, mathtools). The delivered numbering is the unit's own. Assets: no retained span contains a figure environment, a graphics-inclusion command, a PDF-page inclusion, a table or a TikZ picture; no image file is copied into this package, the package declares no asset dependency and no image file is redistributed with it. Licence: version arXiv:2309.05657v3 of this article is distributed under the Creative Commons Attribution 4.0 International licence, as recorded in the arXiv licence metadata for this exact version and shown on the abstract page https://arxiv.org/abs/2309.05657v3. A full rights notice is delivered at licenses/arxiv-2309.05657-CC-BY-4.0.txt. Characters: every retained excerpt is pure ASCII, so the delivered engine, XeLaTeX with its default Latin Modern text font, reports no missing character.
\begin{lemma}
\label{density lower bound}
Let $f(t) = e^{-g(t)}$ be a log-concave probability density on $\R$
having variance $1$ and let $m$ denote its (unique) median. Then 
$$ e^{-g(m)}\geq \frac{e^{-4}}{2}~.$$
\end{lemma}

\begin{proof}
Without loss of generality, we may assume that $m=0$ and $g$ takes
its minimum on $\R^-$. Following the remark at the top of the appendix, we also assume that $g$ is continuous.

  % The intuition is that if the density is too low in $0$, it has to be almost constant on a wide set around $0$ to have total mass $1$. But then the variance will be greater than 1. 
If $g(0)\leq 0$ the result is obvious, so suppose $g(0)>0$.
Since $g$ is convex, by taking its minimum on $\R^-$ it is non
decreasing on $\R^+$.
By continuity, there exists $L>0$ such that
$g(L)=2g(0)$.

From the convexity of $g$ we have that $g'(L)\geq \frac{g(0)}{L}$, and
therefore  for all $t \geq L$, 
\begin{equation}\label{ineq 5}
g(t)\geq g(L)+\frac{g(0)}{L}(t-L)\geq \frac{g(0)}{L}t~.
\end{equation}
Since $\int_{\R} f(x) dx =1$ and $0$ is the median,
\[
  \frac{1}{2}=\int_0^L e^{-g(t)}dt+\int_L^{\infty} e^{-g(t)}dt~.
\]
Using \eqref{ineq 5},
\[
  \int_L^{\infty} e^{-g(t)}dt \leq \int_L^{\infty}
  e^{-g(0)t/L}dt=\frac{L}{g(0)}e^{-g(0)}~.
\]
Because $g$ is non-decreasing on $\R^+$, 
\[
  \int_0^L e^{-g(t)}dt \leq e^{-g(0)}L~,
\]  
leading to
\begin{equation}\label{ineq 6}
\frac{1}{2} \leq \frac{L}{g(0)}e^{-g(0)}+e^{-g(0)}L=e^{-g(0)}L\left(1+\frac{1}{g(0)}\right)~.
\end{equation}
Now we use the fact that the variance equals $1$, that is,
\[
  1=\int_{-\infty}^{+\infty}t^2e^{-g(t)}dt -
  \left(\int_{-\infty}^{\infty}t e^{-g(t)}dt\right)^2~.
\]
Since the difference between the expectation and the median of any
distribution is at most the standard deviation, we have
$|\int_{-\infty}^{\infty}t e^{-g(t)}dt|\le 1$.
Moreover,  since $g$ is increasing on $\R^+$, for all $t \in [0,L]$ we
have $g(t)\leq 2g(0)$, and therefore $1 \geq
\int_0^{\infty}t^2e^{-g(t)}dt -1$
implies
\begin{equation}\label{ineq 7}
2 \geq \int_0^L t^2 e^{-2g(0)}dt = \frac{L^3}{3}e^{-2g(0)}~.
\end{equation}
From \eqref{ineq 6} we have
\[
  e^{-2g(0)}L^3  e^{-g(0)}\left(1+\frac{1}{g(0)}\right)^3 \geq
  \frac{1}{8}~.
\]
Hence, by plugging the inequality into \eqref{ineq 7}, we get
\begin{equation}\label{ineq 8}
e^{-g(0)}\left(1+\frac{1}{g(0)}\right)^3 \geq \frac{1}{48}~.
\end{equation}
Note that the function $h:t \mapsto e^{-t}\left(1+\frac{1}{t}\right)^3$
is non increasing on $\R^+ $.
To conclude, observe that
\begin{itemize}
\item if $g(0)\leq 4.5$, then $e^{-g(0)}\geq \frac{e^{-4}}{2}$.
\item if $g(0)>4.5$, then
$$ h(g(0)) < \frac{1}{48},$$
contradicting \eqref{ineq 8}.
\end{itemize}
\end{proof}

\end{document}
